% !TEX root = ../Grunddatei.tex
\section{Conclusion}

The results shown indicate that the proposed workflow can be used to detect
one-way streets, by combining satellite imagery, street view images and modern
deep learning technologies. The presented application is able to produce
meaningful and accurate results that can be used for further analysis.

However, the current workflow also introduces some limitations and challenges. First, the approach and implementation lean heavily on the existence of available, high-quality image data at a given location. Since streets and their surrounding environment are highly complex frameworks, it can be possible that the image quality is insufficient for an effective segmentation of intersections and thus, the entire analysis. It is also possible for the available satellite images to stem from system that can only offer low-resolution imagery, which would limit the quality and accuracy of the procedure. Another limiting factor is that, as discussed, Panoramax are mostly restricted to French territory. Since the application uses 360-degree images, it is obvious that the proposed workflow does not work in areas with low density of street view panoramas. To address this issue, other services, like Google Street View, could be included, although that would most likely mean including proprietary data sources.

Further development of the approach could include refining the current implementation or model architecture to achieve more accurate results that are less vulnerable to errors. It is also possible to test and expand the workflow for other areas, cities or countries. Another advancement could be to more accurately represent different edge cases and make the application more resilient to the complex nature of road environment analysis.